\documentclass[tikz,border=4pt]{standalone}
\usepackage{tikz}
\usepackage{polymer-paper-preamble}

\begin{document}
\begin{tikzpicture}[
  x=1cm,
  y=1cm,
  line cap=round,
  line join=round,
  font=\sffamily\fontsize{6.2}{7.4}\selectfont,
  panel/.style={draw=cLGray, fill=white, rounded corners=2pt, line width=0.45pt},
  panel title/.style={font=\sffamily\bfseries\fontsize{8.2}{9.8}\selectfont, text=black, anchor=west},
  label/.style={font=\sffamily\fontsize{6.2}{7.4}\selectfont, text=cGray, align=center},
  small label/.style={font=\sffamily\fontsize{5.5}{6.6}\selectfont, text=cGray, align=center},
  axis/.style={draw=cGray, line width=0.42pt, -{Stealth[length=2.5pt,width=1.8pt]}},
  mechanism arrow/.style={draw=cTeal, line width=0.85pt, -{Stealth[length=3.2pt,width=2.2pt]}},
  trap mark/.style={draw=cAmber, line width=0.8pt}
]

% Compact two-block schematic
\draw[panel] (0,0) rectangle (5.75,4.75);
\draw[panel] (5.95,0) rectangle (11.55,4.75);
\node[panel title] at (0.20,4.47) {(a) Trapping during conduction};
\node[panel title] at (6.15,4.47) {(b) Trap-energy distribution};

% Panel A: applied bias cell
\draw[fill=cGray, draw=cGray, line width=0.35pt, rounded corners=0.5pt]
  (0.48,1.05) rectangle (0.72,3.36);
\draw[fill=cGray, draw=cGray, line width=0.35pt, rounded corners=0.5pt]
  (2.34,1.05) rectangle (2.58,3.36);
\draw[fill=cLGray!30, draw=cGray, line width=0.38pt]
  (0.72,1.05) rectangle (2.34,3.36);
\node[small label, text=white, rotate=90] at (0.60,2.20) {electrode};
\node[small label, text=white, rotate=90] at (2.46,2.20) {electrode};

\node[label, text=black, anchor=east] at (1.20,3.74) {applied};
\node[font=\sffamily\bfseries\fontsize{7.0}{8.4}\selectfont, text=cBlue, anchor=west]
  at (1.53,3.74) {$V$};
\draw[draw=cBlue, line width=0.62pt, {Stealth[length=2.5pt,width=1.8pt]}-{Stealth[length=2.5pt,width=1.8pt]}]
  (0.58,3.48) -- (2.48,3.48);
\node[small label, text=black] at (1.53,0.78) {disordered sulfur--polymer film};

% Unlabelled trap sites and a non-ballistic, repeated-trapping walk
\foreach \x/\y in {0.97/1.37,1.66/1.38,1.21/1.76,2.03/1.92,1.50/2.25,0.99/2.64,1.89/2.89}{
  \draw[trap mark] (\x-0.055,\y-0.055) -- (\x+0.055,\y+0.055)
                   (\x-0.055,\y+0.055) -- (\x+0.055,\y-0.055);
}
\draw[mechanism arrow]
  (1.02,3.03)
  .. controls (1.52,3.29) and (2.03,2.91) .. (1.50,2.25)
  .. controls (1.10,1.91) and (1.17,1.57) .. (1.66,1.38)
  .. controls (2.11,1.21) and (2.19,1.70) .. (2.03,1.92)
  .. controls (1.78,2.26) and (1.43,2.04) .. (1.21,1.76);
\draw[draw=cTeal, fill=white, line width=0.65pt] (1.02,3.03) circle (0.09);
\node[small label, text=cTeal, anchor=west, align=left] at (0.78,3.19)
  {sign-agnostic carrier};
\node[small label, text=black] at (1.54,0.50)
  {tortuous walk $\cdot$ repeated trapping};

% Qualitative transient-current inset
\draw[axis] (3.03,1.12) -- (5.27,1.12)
  node[font=\sffamily\fontsize{5.5}{6.6}\selectfont, text=cGray, below left=1pt] {$t$};
\draw[axis] (3.03,1.12) -- (3.03,3.34)
  node[font=\sffamily\fontsize{5.5}{6.6}\selectfont, text=cGray, below right=1pt] {$I$};
\draw[draw=cBlue, line width=1.0pt]
  (3.20,3.10) .. controls (3.50,2.24) and (4.07,1.60) .. (5.10,1.31);
\node[label, text=black, anchor=west, align=left] at (3.16,3.55)
  {absorption current};
\node[font=\sffamily\bfseries\fontsize{6.5}{7.8}\selectfont, text=cBlue,
  fill=white, inner sep=1.2pt] at (4.12,2.30) {$I(t)\propto t^{-n}$};
\node[font=\sffamily\bfseries\fontsize{6.5}{7.8}\selectfont, text=cBrown,
  fill=cAmberSphere!35, rounded corners=1.5pt, inner xsep=4pt, inner ysep=2.5pt]
  at (4.15,0.58) {$I\!\downarrow\ \Rightarrow\ R\!\uparrow$};
\draw[draw=cGray, line width=0.45pt, -{Stealth[length=2.5pt,width=1.8pt]}]
  (2.70,2.18) -- (2.95,2.18);

% Panel B: common energy-distribution axes
\draw[axis] (6.46,1.02) -- (11.08,1.02)
  node[font=\sffamily\fontsize{5.5}{6.6}\selectfont, text=cGray, below left=1pt]
  {trap energy, $E$};
\draw[axis] (6.46,1.02) -- (6.46,3.87)
  node[font=\sffamily\fontsize{5.5}{6.6}\selectfont, text=cGray, below right=1pt]
  {$g(E)$};

% S60: discrete shallow state plus one dominant narrow deep state
\draw[draw=cAmber, line width=0.95pt]
  (6.70,1.02) .. controls (6.82,1.06) and (6.86,1.51) .. (6.95,1.84)
  .. controls (7.04,1.51) and (7.09,1.06) .. (7.22,1.02);
\draw[draw=cAmber, line width=1.05pt]
  (7.45,1.02) .. controls (7.62,1.10) and (7.70,2.67) .. (7.86,3.22)
  .. controls (8.02,2.67) and (8.12,1.10) .. (8.31,1.02);
\node[label, text=cBrown] at (7.48,3.58)
  {$\mathrm{S60}$: discrete shallow + deep};
\node[small label, text=cBrown] at (7.86,2.62) {deep};

% S80: continuous broad landscape
\path[fill=cTeal!22, draw=none]
  (8.48,1.02)
  .. controls (8.84,1.11) and (9.16,2.22) .. (9.67,2.67)
  .. controls (10.18,3.11) and (10.69,2.31) .. (10.96,1.02)
  -- cycle;
\draw[draw=cTeal, line width=1.05pt]
  (8.48,1.02)
  .. controls (8.84,1.11) and (9.16,2.22) .. (9.67,2.67)
  .. controls (10.18,3.11) and (10.69,2.31) .. (10.96,1.02);
\node[label, text=cTeal] at (9.75,3.58)
  {$\mathrm{S80}$: continuous broad};

% Breadth and magnitude are deliberately orthogonal cues
\draw[draw=cTeal, line width=0.68pt,
  {Stealth[length=2.5pt,width=1.8pt]}-{Stealth[length=2.5pt,width=1.8pt]}]
  (8.64,1.47) -- (10.81,1.47);
\node[font=\sffamily\bfseries\fontsize{6.1}{7.3}\selectfont, text=cTeal,
  fill=white, inner sep=1pt] at (9.72,1.47) {$n$ = breadth};
\draw[draw=cGray, line width=0.62pt,
  {Stealth[length=2.5pt,width=1.8pt]}-{Stealth[length=2.5pt,width=1.8pt]}]
  (11.20,1.08) -- (11.20,2.73);
\node[font=\sffamily\fontsize{5.7}{6.8}\selectfont, text=cGray, rotate=90,
  fill=white, inner sep=1pt] at (11.20,1.90)
  {$\rho_{60\mathrm{s}}$ = magnitude};

% Composition-driven evolution
\draw[draw=cAmber, line width=0.70pt, -{Stealth[length=3pt,width=2pt]}]
  (7.78,4.02) -- (9.88,4.02);
\node[small label, text=cBrown, fill=white, inner sep=1pt] at (8.83,4.02)
  {sulfur / disorder $\uparrow$};
\node[label, text=black] at (8.76,0.42)
  {discrete $\longrightarrow$ continuous broad landscape};

\end{tikzpicture}
\end{document}
